\documentclass[10pt,a4paper]{amsart}
\usepackage[margin=19mm]{geometry}
\usepackage{amsmath,amssymb,mathtools}
\usepackage[scr=rsfs,cal=euler]{mathalfa}
\usepackage{xfrac}
\usepackage[unicode,hidelinks,bookmarks=false]{hyperref}
\newcommand{\ie}{i.e.\ }
\newcommand{\cf}{cf.\ }
\newcommand{\resp}{respectively\ }
\newcommand{\Subsection}{\subsection}
\providecommand{\nobreakdash}{\nobreak\hbox{-}}
\setlength{\emergencystretch}{2em}
\multlinegap0pt
\usepackage{amssymb}
\usepackage{xcolor}
\usepackage{tikz}
\newtheorem{theorem}{Theorem}
\newcommand{\mustar}{\mu^{*}}
\title{Sharp bounds between the saturation number\\ and the harmonic index}
\author{Chakshu Gupta}
\date{Source: arXiv:2606.15761v3 (CC BY 4.0)}
\begin{document}
\maketitle

\noindent\emph{Source and licence.} Sharp bounds between the saturation number\\ and the harmonic index. Version arXiv:2606.15761v3 (CC BY 4.0), distributed by arXiv under the Creative Commons Attribution 4.0 International licence, http://creativecommons.org/licenses/by/4.0/, as recorded in the arXiv licence metadata for this exact version and shown on the abstract page https://arxiv.org/abs/2606.15761v3. Full licence notice: \texttt{licenses/arxiv-2606.15761-CC-BY-4.0.txt}.\par\smallskip

\noindent\textit{Editorial scope.} Counted unit: the article's theorem thm:treebound (source0.tex lines 275--278). The article's complete original proof of it is delivered with it (source0.tex lines 280--371, one \texttt{proof} environment of 92 lines). No mathematics is written, repaired or re-derived: the statement and the proof are the article's own lines, in the article's own notation, and no retained line is substituted or re-flowed.  Retained uncounted as the source-local context the proof needs: lines 225--237.  Nothing was omitted from any retained span, so the package carries no omission record.  No retained line contains a citation, so the wrapper carries no bibliography and no bibliography entry.  No retained line contains a figure environment, an image-inclusion command or drawing code, so no image is delivered and the package declares no asset dependency.  Editorial formatting only, in addition: an AMS \texttt{amsart} preamble that restates the article's own declarations and macros used by the retained lines, namely the environments \texttt{theorem} on separate counters, together with the standard packages those lines need.  The retained environments print in this document as Theorem 1 in order of appearance, whereas the article prints the same items with the numbering of its own full document.  The complete native source of the article as distributed in the arXiv e-print for this exact version is bundled unchanged as \texttt{sources/202886/source0.tex}.
\begin{theorem}[Subdivided star]\label{thm:tree}
For the subdivided star~$S_k$ ($k \ge 1$),
\[
  \mustar(S_k) = k
  \qquad\text{and}\qquad
  H(S_k) = \frac{2k}{k+2} + \frac{2k}{3}.
\]
Consequently $\mustar(S_k) \le H(S_k)$ if and only if $k \le 4$, with
equality at $k = 4$ ($S_4$, nine vertices), and
$\mustar(S_k)/H(S_k) \to 3/2$ as $k \to \infty$.  In particular, the
conjecture fails for trees, and hence for triangle-free graphs and for
bipartite graphs.
\end{theorem}

\begin{theorem}[Tree bound]\label{thm:treebound}
Every nontrivial tree~$T$ satisfies $\mustar(T) < \tfrac{3}{2} H(T)$, and the
constant~$3/2$ is best possible.
\end{theorem}

\begin{proof}
Since every maximal matching is a matching, $\mustar(T) \le \nu(T)$, the
matching number.  It therefore suffices to prove the stronger inequality
$2\nu(T) < 3H(T)$ for every nontrivial tree~$T$, by strong
induction on $|V(T)|$.

\emph{Base case.}  If $T$ has no path on four vertices, it is a star
$K_{1,k}$ with $k \ge 1$, a centre joined to $k$ leaves.  Every edge meets
the centre, so a matching uses at most one of them and $\nu(T) = 1$.  Each
edge joins the centre of degree~$k$ to a leaf of degree~$1$, so
$H(T) = k \cdot 2/(k+1) = 2k/(k+1)$.  Hence
$3H(T) - 2\nu(T) = (4k-2)/(k+1) > 0$.

\emph{Inductive step.}  Otherwise $T$ has a path on four vertices.
Fix a longest one, $v_0 v_1 \cdots v_L$ with $L \ge 3$, and write $s = v_1$
and $p = v_2$.  By
maximality of the path, every neighbour of $s$ other than $p$ is a leaf,
since a non-leaf neighbour would extend the path beyond $v_0$.  Let
these leaves be $\ell_1, \ldots, \ell_t$, so $v_0$ is among them, $t \ge 1$, and
$d(s) = t+1$.  Since $p$ is adjacent to both $s$ and $v_3$, its degree
$D := d(p)$ is an integer at least~$2$.  Let
$T^- = T - \{s, \ell_1, \ldots, \ell_t\}$.
Each $\ell_i$ has $s$ as its only neighbour, and $s$ is otherwise adjacent
only to $p$, so $sp$ is the only edge joining the removed vertices to the
rest of $T$.  Hence $T^-$ is a tree containing $p$ and $v_3$, nontrivial and
smaller than~$T$.

The leaf $\ell_1$ has $s$ as its only neighbour, so some maximum matching
contains $s\ell_1$.  Deleting $s$ and $\ell_1$ then isolates
$\ell_2, \ldots, \ell_t$, so $\nu(T) = \nu(T^-) + 1$.  For the harmonic
index, the only surviving vertex whose degree changes is $p$, falling from
$D$ to $D-1$.  The deleted edges are the $t$ leaf edges $s\ell_i$, each of
weight $2/(t+2)$, and the edge $sp$, of weight $2/(t+1+D)$.  Write
$z_1, \ldots, z_{D-1}$ for the neighbours of $p$ other than $s$.  Each
remaining edge $pz_j$ rises in weight from $2/(D+d(z_j))$ to
$2/(D-1+d(z_j))$.  Let $\Lambda$ be the total of these increases:
\[
  \Lambda = \sum_{j=1}^{D-1}
      \left( \frac{2}{D-1+d(z_j)} - \frac{2}{D+d(z_j)} \right).
\]
Then
\[
  H(T) - H(T^-) = \frac{2t}{t+2} + \frac{2}{t+1+D} - \Lambda.
\]
Each summand of $\Lambda$ combines into a single fraction, which $d(z_j) \ge 1$
bounds:
\begin{align*}
  \frac{2}{D-1+d(z_j)} - \frac{2}{D+d(z_j)}
    &= \frac{2}{(D-1+d(z_j))(D+d(z_j))} \\
    &\le \frac{2}{D(D+1)}.
\end{align*}
Summing the $D-1$ terms gives $\Lambda \le 2(D-1)/[D(D+1)]$, so
\[
  H(T) - H(T^-) \ge \frac{2t}{t+2} + \frac{2}{t+1+D}
                     - \frac{2(D-1)}{D(D+1)}.
\]
It remains to bound the right-hand side.  The values $t$ and $D$ depend
on~$T$, but satisfy $t \ge 1$ and $D \ge 2$, so it suffices to show the
right-hand side exceeds $\tfrac{2}{3}$ for all such integers.  Since
$\tfrac{2t}{t+2} - \tfrac{2}{3} = \tfrac{4(t-1)}{3(t+2)}$, this reduces to
\[
  \frac{4(t-1)}{3(t+2)} + \frac{2}{t+1+D}
    - \frac{2(D-1)}{D(D+1)} > 0.
\]
For $t = 1$ the first term vanishes, and
\[
  \frac{2}{D+2} - \frac{2(D-1)}{D(D+1)} = \frac{4}{D(D+1)(D+2)} > 0,
\]
using $D(D+1) - (D-1)(D+2) = 2$.  For $t \ge 2$,
\[
  \frac{4(t-1)}{3(t+2)} \ge \frac{1}{3}
  \qquad\text{and}\qquad
  \frac{2(D-1)}{D(D+1)} \le \frac{1}{3},
\]
the first with equality at $t = 2$, the second because
$(D-2)(D-3) \ge 0$ for every integer $D \ge 2$.  The middle term
$2/(t+1+D)$ is positive, so the sum is at least
$\tfrac{1}{3} + 2/(t+1+D) - \tfrac{1}{3} > 0$.  Thus
$H(T) - H(T^-) > \tfrac{2}{3}$.

Combining the two drops,
\begin{align*}
  3H(T) - 2\nu(T)
    &= \bigl(3H(T^-) - 2\nu(T^-)\bigr)
      + 3\bigl(H(T) - H(T^-)\bigr) - 2\bigl(\nu(T) - \nu(T^-)\bigr) \\
    &> \bigl(3H(T^-) - 2\nu(T^-)\bigr) + 3 \cdot \tfrac{2}{3} - 2,
\end{align*}
which equals $3H(T^-) - 2\nu(T^-) > 0$ by the induction hypothesis.  Hence
$2\nu(T) < 3H(T)$, and so $\mustar(T) \le \nu(T) < \tfrac{3}{2} H(T)$.
By Theorem~\ref{thm:tree}, $\mustar(S_k)/H(S_k) \to \tfrac{3}{2}$ as
$k \to \infty$, so no smaller constant suffices.
\end{proof}

\end{document}
